\documentclass{article}
\usepackage[utf8]{inputenc}

\usepackage[left=2.5cm, right=2.5cm, top=3cm, bottom=2cm]{geometry}

\usepackage{tikz,amsmath, amssymb,bm,color}
\usetikzlibrary{shapes,arrows}

% needed for BB
\usetikzlibrary{calc}
\usepackage{pgfplots}
\pgfplotsset{compat=1.16}
\usepackage{circuitikz}
\usetikzlibrary{patterns}

\usepackage[ngerman]{babel}
\usepackage{hyperref}
\usepackage[ngerman]{cleveref}

\renewcommand{\vec}{\boldsymbol}

\title{Verdrehtes Elektrisches Feld}
\author{ Joachim Schwardt und Tim Pokart}
\date{\today}

\begin{document}

\maketitle

Wir betrachten den in \cref{fig:sketch} dargestellten Aufbau, bei dem zwei ineinandergestapelte Hohlzylinder mit Querschnittsfläche $A$ jeweils gegenläufig vom Strom $I$ durchflossen werden. 
Der innere hat dabei den Außenradius $a$ und der äußere den Innenradius $b$.
\begin{figure}[htb]
    \centering
    \begin{tikzpicture}

\draw[blue,pattern=north east lines,pattern color=blue] (1, 0) -- (1.5,0) -- (1.5,-2.5) node[midway,right]{A} --  (1,-2.5) --cycle;

\draw[-latex]  (1.05,0.45) arc (45:135:1.5 and 0.5) node[above, midway] {$I$};

\draw[dotted] (0,0) -- (0,-2.5);


\draw  (0,0) ellipse (1 and 0.25);
\draw  (0,0) ellipse (1.5 and 0.5);

\draw[dashed]  (-1,0) edge (-1,-2.5);
\draw  (-1.5,0) edge (-1.5,-2.5);
\draw[dashed]  (1,0) edge (1,-2.5);
\draw  (1.5,0) edge (1.5,-2.5);

\draw[dashed]  (0,-2.5) ellipse (1 and 0.25);
\draw  (-1.5,-2.5) arc (180:360:1.5 and 0.5);
\draw[dashed]  (1.5,-2.5) arc (0:180:1.5 and 0.5);




\draw[|-|] (-1.75,0) -- (-1.75,-2.5) node[midway, left] {$h$};

\draw[|-|] (-1.5,-1.25) -- (0,-1.25) node[midway, above] {$a$};

\begin{scope}[yshift=-120.0]
\draw[latex-]  (1.05,0.45) arc (45:135:1.5 and 0.5) node[below, midway] {$I$};
\draw[blue,pattern=north east lines,pattern color=blue] (2, 0) -- (2.5,0) -- (2.5,-2.5) node[midway,right]{A} --  (2,-2.5) --cycle;
\draw[dotted] (0,0) -- (0,-2.5);

\draw  (0,0) ellipse (2 and 0.75);
\draw  (0,0) ellipse (2.5 and 1.0);

\draw[dashed]  (-2,0) edge (-2,-2.5);
\draw  (-2.5,0) edge (-2.5,-2.5);
\draw[dashed]  (2,0) edge (2,-2.5);
\draw  (2.5,0) edge (2.5,-2.5);

\draw[dashed]  (0,-2.5) ellipse (2 and 0.75);
\draw  (-2.5,-2.5) arc (180:360:2.5 and 1.0);
\draw[dashed]  (2.5,-2.5) arc (0:180:2.5 and 1.0);

\draw[|-|] (-2.75,0) -- (-2.75,-2.5) node[midway, left] {$h$};

\draw[|-|] (-2.0,-1.35) -- (0,-1.35) node[midway, above] {$b$};


\draw[dotted,very thin]  (0,0) ellipse (1 and 0.25);
\draw[dotted,very thin]  (0,0) ellipse (1.5 and 0.5);

\draw[dotted,very thin]  (-1,0) edge (-1,-2.5);
\draw[dotted,very thin]  (-1.5,0) edge (-1.5,-2.5);
\draw[dotted,very thin]  (1,0) edge (1,-2.5);
\draw[dotted,very thin]  (1.5,0) edge (1.5,-2.5);

\draw[dotted,very thin]  (0,-2.5) ellipse (1 and 0.25);
\draw[dotted,very thin]  (-1.5,-2.5) arc (180:360:1.5 and 0.5);
\draw[dotted,very thin]  (1.5,-2.5) arc (0:180:1.5 and 0.5);

\end{scope}


\end{tikzpicture}
    \caption{Skizze der Anordnung}
    \label{fig:sketch}
\end{figure}
Gesucht ist das Elektrische Feld im Zwischenraum.
% Dafür verwenden wir lediglich die Maxwell-Gleichungen
% \begin{align*}
%     \nabla\cdot \vec{D} &= \rho, &\nabla\cdot \vec{B} &= 0, & \nabla \times \vec{E} + \frac{\partial \vec{B}}{\partial t} &= 0, & \nabla \times \vec{H} - \frac{\partial \vec{D}}{\partial t} &= \vec{j},
% \end{align*}
% das Elektrische Potential $\nabla \varphi = -\vec{E}$ und die sich daraus ergebende Poisson-Gleichung $\Delta \varphi = -\frac{\rho}{\epsilon_0}$, das ohmsche Gesetz $\vec{j} = \kappa \vec{E}$ mit der Leitfähigkeit $\kappa$, sowie die Kontinuitätsbedingung des Elektrischen Feldes an Grenzübergängen zwischen Stoffen $\vec{n} \times (\vec{E}_2 - \vec{E}_1) = 0$.
Dafür verwenden wir die Maxwell-Gleichungen in der Form
\begin{align*}
    \nabla\cdot \vec{E} &= \frac{\rho }{\epsilon_0}, & \nabla \times \vec{E} &= 0,
\end{align*}
wir nehmen also zeitunabhängige Felder an. 
Da die Stromdichten zeitunabhängig sind, muss das auch für die Felder gelten -- auch wenn das Endergebnis einer ähnlichen Symmetrieüberlegung widerspricht.
Das elektrische Feld lässt sich aufgrund der Rotationsfreiheit über das elektrostatische Potential als $\vec{E} = -\nabla \varphi$ ausdrücken.
Zusammen mit der Tatsache, dass aufgrund fehlender freier Ladungen $\rho = 0$ ist, ergibt sich daraus die Poisson-Gleichung $\Delta \varphi = -\frac{\rho}{\epsilon_0} \equiv 0$.
Wir nehmen zudem noch das ohmsche Gesetz in seiner einfachsten Form $\vec{j} = \kappa \vec{E}$ an, wobei $\kappa$ die Leitfähigkeit ist.
Zuletzt verwenden wir noch die Stetigkeit der Parallelkomponente des elektrischen Feldes an Grenzübergängen, welche mit dem Normalenvektor $\vec{n}$ der Grenzfläche als $\vec{n} \times (\vec{E}_2 - \vec{E}_1) = 0$ formuliert werden kann.
% Wir arbeiten in den Zylinderkoordinaten $(r, \phi, z)$ und beschränken uns der Einfachheit halber auf $z=0$. 
% Für unsere Betrachtung reicht es sich auf die Koordinaten $r = a$ und $r = b$ zu beschränken, da diese die Anschlussbedingungen für die spätere Lösung der Poisson-Gleichung stellen.
Wir arbeiten in den Zylinderkoordinaten $(r, \phi, z)$ und beschränken uns im Folgenden auf $z=0$. 
Für die Lösung der Poisson-Gleichung benötigen wir Randbedingungen, die in diesem Fall den Parallelkomponenten des elektrischen Feldes bei $r = a$ und $r = b$ entsprechen.
% Zunächst bestimmen wir die Stromdichte 
% \begin{align*}
%     \vec{j} = \pm\vec{e}_\phi \frac{I}{A} \Theta\left(\frac{h}{2} - |z|\right)
% \end{align*}
% in den beiden Leiter an eben diesen Stellen und daraus das Elektrische Feld 
% \begin{align*}
%     \vec{E}(r = a, \phi, z) &= \vec{e}_\phi \frac{I}{\kappa A} \Theta\left(|z| - \frac{h}{2}\right), &\vec{E}(r = b, \phi, z) &=-\vec{e}_\phi \frac{I}{\kappa A} \Theta\left(\frac{h}{2} - |z|\right).
% \end{align*}
Die Stromdichte ist gegeben durch
\begin{align*}
    \vec{j}(z=0) = \pm\vec{e}_\phi \frac{I}{A}
\end{align*}
und entsprechend folgt dann
\begin{align*}
    \vec{E}(r = a, \phi, z=0) &= \vec{e}_\phi \frac{I}{\kappa A}, &\vec{E}(r = b, \phi, z=0) &=-\vec{e}_\phi \frac{I}{\kappa A}.
\end{align*}
% Diese beiden Werte sind die Neumann-Randbedingung für die Poisson-Gleichung $\Delta \varphi = 0$, deren Lösung durch diese bis auf einen konstanten Term eindeutig bestimmt wird.
Wir wählen den Ansatz $\varphi(r, \phi, z = 0) = \phi f(r)$\footnote{Ohne die $\phi$-Abhängigkeit können später die Randbedingungen nicht erfüllt werden.} für den Zwischenraum der beiden Leiter in der $x$-$y$-Ebene.
Mit dem Laplace-Operator in Zylinderkoordinaten erhalten wir
\begin{align*}
    \Delta \varphi = \frac{1}{r} \partial_r \left(r \partial_r \varphi\right) + \frac{1}{r^2} \partial_\phi^2 \varphi + \partial_z^2 \varphi = 0.
\end{align*}
Daraus ergibt sich eine Differentialgleichung für $f(r)$ 
\begin{align*}
    \partial_r (r f'(r)) = 0.
\end{align*}
Diese wird durch $f(r) = B \ln r + C$ gelöst.
Damit ergibt sich das elektrische Feld 
\begin{align*}
    \vec{E} &= -\nabla \varphi = -\vec{e}_r \phi \frac{B}{r} - \vec{e}_\phi \frac{B \ln r + C}{r}.
\end{align*}
Dieses kann nun an die beiden Randbedingungen angepasst werden, wobei der Normalenvektor an die Ebene $\vec{n} = \vec{e}_r$ ist.
Zusammen mit $\vec{n} \times (\vec{E}_2 - \vec{E}_1) = 0$ ergeben sich die Gleichungen
\begin{align*}
    \vec{e}_z \left(\frac{I}{\kappa A} + \frac{B \ln a + C}{a}\right) &= 0, \\
    \vec{e}_z \left(-\frac{I}{\kappa A} + \frac{B \ln b + C}{b}\right) &= 0.
\end{align*}
Diese Gleichungen sind äquivalent zu
\begin{align*}
    C &= -\frac{Ia}{\kappa A} - B\ln a,\\
    0 &= -\frac{Ib}{\kappa A} + B\ln \frac{b}{a} - \frac{Ia}{\kappa A}
\end{align*}
und liefern letztlich
\begin{align*}
    B &= \frac{I}{\kappa A} \frac{a+b}{\ln \frac{b}{a}},\\
    C &= -\frac{I}{\kappa A} \left( a + \frac{a+b}{\ln \frac{b}{a}} \ln a \right).
\end{align*}
% Dieses wird durch
% \begin{align*}
%     B &= \frac{I}{\kappa A} \frac{a + b}{\ln \frac{a}{b}}, & C&=\frac{I}{\kappa A} \left(a - \frac{a - b}{\ln \frac{a}{b}} \ln a\right)
% \end{align*}
% gelöst.
Damit hat unser Elektrisches Feld schlussendlich eine winkelabhängige Komponente in radialer Richtung, obwohl das Problem rotationssymmetrisch ist.
Insbesondere ist $\vec{E}(\phi=0) \neq \vec{E}(\phi=2\pi)$.\\
SANITY CHECK: Rotation gleich Null?
\begin{align*}
    \nabla \times \vec{E} &= \left( \vec{e}_r\partial_r + \frac{1}{r}\vec{e}_\phi \partial_\phi + \vec{e}_z\partial_z \right) \times \left( -\vec{e}_r\phi \frac{B}{r} - \vec{e}_\phi \frac{B\ln r + C}{r} \right) = \\
    &= \vec{e}_z \left( -\partial_r \left[ \frac{B\ln r + C}{r} \right] + \frac{B}{r^2} \right) - \vec{e}_z \frac{B\ln r + C}{r^2} = 0.
\end{align*}

\end{document}
